\documentclass[exercice,a4]{thedocument}%
\usepackage{texmaths-tikz}%
\usepackage{pstricks}%
\usepackage{pst-bezier}%
\usepackage{pstricks-add}%
\usepackage{pst-eucl}%
\usepackage{pst-func}%
\usepackage{pst-tree}%
\usepackage{pst-circ}%
\usepackage{tikz}%
\usepackage{pgfplots}%
\usepackage{tkz-tab}%
\usepackage{tikzpagenodes}%
\usepackage{tabularray}%
\usepackage{tabularx}%
\usepackage{multirow}%
\usepackage{multicol}%
\usepackage{scratch3}%
\usepackage{enumitem}%
\usepackage[np]{numprint}%
\usepackage{pas-tableur}%
\usepackage{listings}%
\usepackage{graphicx}%
\usepackage{adjustbox}%
\usepackage{needspace}%
\usepackage{xcolor}%
\usepackage{eurosym}%
\usepackage{diagbox}%
\usepackage{colortbl}%
\usepackage{multido}%
\usepackage{icomma}%
\usepackage[french]{babel}%
\graphicspath{{../brevet/2026/}}%
\startdatakeys
    \source{26GENMATAG1}
    \BO{2026}
    \cycle{4}
    \competencesDuSocle{CA,RA,CO}
    \objectifsApprentissage{C4A1aa,C4A2ad,C4B2cb,C4D2aa,C4D3ac}
    \motsCles{Pythagore,trigonométrie,angle,réciproque de Thalès,homothétie}
\enddatakeys
\definecolor{gris}{gray}{0.5}%
\providecommand{\wideparen}[1]{\overset{\frown}{#1}}%
\begin{document}%
\providecommand{\degres}{^\circ}%
\begin{minipage}{0.45\linewidth}
	Sur la figure ci-contre, qui n'est pas représentée en vraie grandeur :
	\begin{itemize}
		\item A est le point d'intersection des droites (BE) et (TM);
		\item le triangle ABM est rectangle en A ;
		\item AT = 2,2 m; MT = 7,7 m;
		\item AE = 1,92m et MB = 7,3 m.
	\end{itemize}
\end{minipage}\hfill~
%Début tikz
\begin{tikzpicture}[baseline={(current bounding box.center)}]
	\draw[fill=gray!30] (0,0) rectangle (0.3,0.3) ;
	\draw[line width=1.2pt]
	(0,5)    node[above](B){B}--
	(-2.3,0) node[left] {T}--
	(0,-2)   node[below]{E}--
	(5.4,0)  node[right]{M}--cycle
	(-2.3,0)--(5.4,0)
	(0,-2)--(0,5)
	(0,0)    node[below left] {A};
\end{tikzpicture}
%Fin tikz
%%%Début PSTricks
%\begin{minipage}{0.47\linewidth}
%	\psset{unit=0.8cm}
%	\begin{pspicture}(0,-2.3)(6.2,5)
%		\pspolygon(2.3,5.6)(0.2,0)(2.3,-1.8)(7.6,0)(2.3,5.6)(2.3,-1.8)(7.6,0)(0.2,0)%BTEBMT
%		\psframe(2.3,0)(2.55,0.25)
%		\uput[u](2.3,5.6){B} \uput[l](0.2,0){T} \uput[d](2.3,-1.8){E} \uput[dl](2.3,0){A} \uput[r](7.6,0){M}
%	\end{pspicture}
%\end{minipage}
%%%Fin PSTricks
\medskip

\begin{enumerate}
	\item Justifier que AM = 5,5~m.

	\item Montrer que la longueur AB est égale à 4,8~m.

	\item Calculer la mesure de l'angle $\widehat{\text{ABM}}$. Donner la valeur arrondie au degré.

	\item Montrer que les droites (TE) et (BM) sont parallèles.

	\textit{Détailler la réponse en précisant la démarche.}

	\item Le triangle ABM est l'image du triangle AET par une homothétie de centre A.

	Choisir le rapport de cette homothétie parmi les propositions suivantes. Aucune justification n'est demandée.

	\begin{center}
		\begin{tabular}{|c|c|c|}
			\hline
			\textbf{Proposition 1} & \textbf{Proposition 2} & \textbf{Proposition 3} \\
			\hline
			2,5 & $-0,4$ & $-2,5$ \\
			\hline
		\end{tabular}
	\end{center}
\end{enumerate}

\needspace{15\baselineskip}
\medskip
\showthesource
\end{document}
